\section{Results}
\label{sec:results}

\subsection{Date uncertainty and context uncertainty differ}

Table~\ref{tab:mechanism} separates parsed questions from contexts containing modeled evidence.
The validation adapter produces 481 claims, of which 221 have explicit end dates.
These claims occupy 1.46\% of the 32,976 indexed source units.
Among 375 parsed questions with modeled evidence in the possible context, 306 have an unchanged-context certificate.
Thus 81.60\% of these contexts need no finer modeled dates to fix their selected evidence.
The other 69 return different contexts under the matched date completions.
\begin{table}[t]
\centering\small
\setlength{\tabcolsep}{4pt}
\begin{tabular}{@{}lrr@{}}
\toprule
Context diagnostic & Pilot & Validation\\
\midrule
All questions & 176 & 2,348 \\
Parsed windows & 173 & 2,296 \\
Modeled context, parsed & 32 & 375 \\
Certified within this group & 23 & 306 \\
Unstable within this group & 9 & 69 \\
Parsed all-fallback context & 141 & 1,921 \\
\bottomrule
\end{tabular}
\caption{Exact context decisions under the fixed fallback policy. Modeled-context rows require at least one modeled claim in the possible top-five context.}
\label{tab:mechanism}
\end{table}


The certificate directs attention to the evidence actually selected.
With modeled evidence anywhere in the first twenty lexical candidates, 987 of 1,056 parsed questions have stable contexts.
Lower-ranked uncertainty can therefore remain unresolved.
All 1,921 parsed all-fallback contexts also remain fixed under the declared policy.
We report them separately because their dates remain unmodeled.

\paragraph{The context budget controls exposure to uncertain dates.}
We retain the same 375-question group and vary only the number of selected source units.
Table~\ref{tab:budget} reports the exact stable-prefix decisions.
Stability falls from 92.27\% at one item to 79.73\% at twenty items.
Most of the decrease occurs within the first five positions.
The diagnostic lets a system distinguish a short stable context from a larger context requiring refinement.
It measures selection before source-text truncation.
\begin{table}[t]
\centering\small
\setlength{\tabcolsep}{4pt}
\begin{tabular}{@{}lrrrrr@{}}
\toprule
Budget $k$ & 1 & 3 & 5 & 10 & 20 \\
\midrule
Stable (\%) & 92.27 & 85.33 & 81.60 & 81.07 & 79.73 \\
\bottomrule
\end{tabular}
\caption{Budget sensitivity on the fixed group of 375 parsed validation questions with modeled evidence at $k=5$. The unit-level certificate precedes source-text truncation.}
\label{tab:budget}
\end{table}


\subsection{Replaying an unstable selection}

For each of the 69 unstable validation contexts, the constructive procedure selects an unresolved claim.
One assignment supports that claim within the query; another excludes it throughout the query.
The replay records include all event dates, source offsets, and selected identifiers.
They connect the abstract certificate to concrete reader inputs.
All 138 constructed worlds replay exactly and produce different rendered contexts within each pair.
In 47 cases, the excluding world retires the selected claim by the query start.
In 22 cases, it moves the claim's start beyond the query end.
Sixty-seven contexts have one unresolved selected claim; two have two.
These are the support decisions that any stabilizing date refinement must settle.

\paragraph{A date boundary changes the evidence.}
Ben Gummer's pilot question asks about May 2017.
The source dates his Ipswich parliamentary tenure to 2010--2017.
Earliest completion chooses January 1, 2017 and removes the tenure excerpt.
Latest completion chooses December 31 and retains it.
The vacated position receives an unrelated heading from the same article.
The certificate traces this difference to the explicit end event.
Refining its bounds entirely after the query start settles inclusion.
An upper end bound at or before that start settles exclusion.
Additional precision therefore needs to resolve this boundary; an exact departure day is unnecessary.


\subsection{Ranking controls separate evidence quality from stability}

Table~\ref{tab:matched} reports exact shown-span coverage under matched source units.
Possible support and the interval control coincide in this run.
The source adapter supplies explicit endpoints and no cross-claim replacement pairs.
Guaranteed support gains 0.09 points of span recall over the interval control.
Its paired interval is $[-0.09,0.26]$ points.
\begin{table}[t]
\centering\small
\setlength{\tabcolsep}{4pt}
\begin{tabular}{@{}lrr@{}}
\toprule
Policy & Span recall & All strings\\
\midrule
No temporal filter & 38.49 & 37.56 \\
Earliest completion & 39.20 & 38.25 \\
Midpoint completion & 39.07 & 38.12 \\
Latest completion & 39.03 & 38.07 \\
Interval control & 39.07 & 38.12 \\
Possible support & 39.07 & 38.12 \\
Guaranteed support & 39.15 & 38.20 \\
Latest-year ranking & 53.51 & 52.51 \\
\bottomrule
\end{tabular}
\caption{Validation evidence coverage on 2,348 questions. Values are percentages. Every policy shares the same source units and context budget. Latest-year ranking reorders the lexical candidates.}
\label{tab:matched}
\end{table}


Latest-year reranking reaches 53.51\% span recall, compared with 39.07\% for the interval control.
The paired difference is 14.44 points, with interval $[12.05,16.86]$.
This result identifies ranking as the main evidence-quality lever in this source adapter.
The context theorem applies to either fixed ranking.
The reported certificate uses the common lexical ranking so temporal policies share identical relevance decisions.

\subsection{Answer generation under matched contexts}
The fixed 24-question reader sample contains 192 method conditions and 52 unique generations.
All outputs pass the frozen JSON parser.
Each method achieves 29.17\% exact answer-set accuracy and 31.94\% strict answer-set F1.
The temporal filters and date completions produce identical answer sets on this sample.
Latest-year ranking changes fourteen answers while retaining the same aggregate strict score.
The nine-question completion-sensitive diagnostic uses 28 unique generations, including six reused from the primary sample.
Earliest and latest date completion change the answer set on four questions.
Possible and guaranteed support also change four answer sets.
Every method has 33.33\% strict answer-set F1 on this selected diagnostic.
These cases show that source-date choices can propagate through context selection into generated answers.



\subsection{Complete answers require complete candidate sets}

Annual controls expose a separate selection loss.
On the full TempLAMA split, latest-value selection achieves 99.89\% first-result accuracy but only 86.37\% answer recall.
Latest-batch selection retains the same first-result accuracy and raises answer recall to 99.98\%.
Selecting only one value suppresses 13.63\% of accepted annual answers before the final context budget.
Its key-cluster confidence interval is $[13.06,14.21]$ points.
The matched Graphiti and exact direct-first policies produce identical endpoints on all 45,866 source observations.
These controls establish the exact-date reference and preserve the benchmark's full answer scope.

\subsection{Independent execution checks}

Enumeration checks the compiler directly against its possible-world definition.
All 253,050 support comparisons and 328,747 context-stability comparisons agree.
Explicit-end verification adds 39,600 support comparisons and 29,700 stability comparisons.
An independent check of the constructive extension validates 1,300 counterworld pairs across 12,925 enumerated worlds.
A separate small-instance check verifies both directions of the Set Cover construction.
These finite correctness checks accompany the proofs and remain separate from corpus measurements.
